\documentclass[10pt]{article}
\usepackage[margin=20mm]{geometry}
\usepackage{amsmath,amssymb,mathtools,fontspec,unicode-math,xurl,hyperref}
\IfFontExistsTF{DejaVu Serif}{\setmainfont{DejaVu Serif}}{\setmainfont{Latin Modern Roman}}
\setmathfont{latinmodern-math.otf}
\setlength{\parindent}{0pt}\setlength{\parskip}{6pt}\sloppy
\title{Initial scientific review of
Generalizing Lusztig’s total positivity}
\author{Paper identifier: M25-0048}
\date{}
\begin{document}
\maketitle
Abstract: This review concerns the assigned publisher version of O. Guichard and A. Wienhard, Inventiones mathematicae 239 (2025), 707–799. One Major finding concerns the representation-theoretic statements used to justify classification and uniqueness of cones. Two Minor findings concern a general restricted-root multiplicity assertion and a globally quantified BCH estimate. Explicit counterexamples and counterevidence are given. These findings do not establish that the four-family classification or the main positivity theorems are false.\par
\section*{Version and assessment}
The reviewed object is the 93-page publisher PDF of [GW], DOI 10.1007/s00222-024-01303-y, published online on 10 January 2025. Printed page numbers equal PDF page numbers plus 706. The independently read Benoist manuscript is identified separately below.\par
The review followed the positivity definitions, representation criterion, classification, explicit invariant cones, open and closed unipotent semigroups, parameter maps, tangent cone, and selected flag-variety arguments. Independent calculations include exact weight and stabilizer computations, a symbolic nilpotent matrix logarithm, an exceptional Weyl-group calculation, and orthogonal rank-two transition equations. No author program was executed.\par
Assessment: one Major finding M01, comprising two related defects in the same representation-theoretic foundation, and two Minor findings M02–M03. The Major label refers to verified invalid inputs in a central proof chain. It does not assert a counterexample to the final classification. The coverage limits in Section 6 (\#sec:limits) are part of this assessment.\par
\section*{M01—The representation criterion and its cone consequence are overstated}
\subsection*{Location and dependency}
Section 3.2, printed p. 725 (PDF p. 19), defines evenness by \(\lambda(H_\alpha)\in2\mathbb Z\) for every weight \(\lambda\) and root \(\alpha\), or \(R(V)\subset2P\), and then asserts that for an irreducible representation it suffices to test the highest weight. The order defining that highest weight is also specified there. Proposition 3.6(1) uses these notions to characterize invariant acute cones; Proposition 3.6(2) identifies the whole interior, up to sign, with vectors fixed by some maximal compact subgroup.\par
These statements are explicit inputs to Proposition 3.7 and the proof of Theorem 3.4 on pp. 726–727 (PDF pp. 20–21). Proposition 3.6(2) is also cited for uniqueness of cones in Remark 3.2, p. 724, and Section 5.1, p. 731. The problem is therefore mathematical, rather than bibliographical.\par
\subsection*{The all-weights criterion fails in an actual classified example}
Equation reference: eq:weights. Let \(S=\operatorname{SL}_3(\mathbb R)\) act on \[V=\operatorname{Sym}^2(\mathbb R^3)=\{M=M^{\mathsf T}\in M_3(\mathbb R)\},\qquad g\cdot M=gMg^{\mathsf T}.\] The positive semidefinite cone is nonzero, closed, convex, pointed, and invariant. Its interior is the positive definite cone. For \(H_1=\operatorname{diag}(1,-1,0)\) and \(H_2=\operatorname{diag}(0,1,-1)\), the weights give \[
\begin{array}{c|rrrrrr}
\lambda&2e_1&2e_2&2e_3&e_1+e_2&e_1+e_3&e_2+e_3\\ \hline
\lambda(H_1)&2&-2&0&0&1&-1\\
\lambda(H_2)&0&2&-2&1&-1&0 .
\end{array}\] The usual highest weight \(2e_1=2\omega_1\) is even on both simple coroots, but three other weights have odd evaluations. Thus the asserted equivalence and Proposition 3.6(1), with the printed all-weights definition, are false.\par
The hypotheses are checkable directly. The six distinct Cartan weights have one-dimensional spaces. Root operators replace one monomial factor \(e_j\) by \(e_i\), connecting all six weight spaces. Cartan spectral projections followed by these operators show that every nonzero invariant subspace is all of \(V\). This proves real irreducibility. The element \(\operatorname{diag}(t,1,t^{-1})\), \(t>1\), has unique largest eigenvalue \(t^2\) on \(V\), with one-dimensional eigenspace, so the representation is proximal in the usual sense.\par
The printed order \(\lambda\leq\mu\) iff \(\lambda(H_\alpha)\leq\mu(H_\alpha)\) for every positive root is another problem: [eq:weights] (\#eq:weights) has no greatest element for this order. Such an element would need both simple-coroot coordinates at least \(2\), and none does. Consequently the claim that every irreducible representation admits a highest weight for that order is false. The usual root-dominance order must be used.\par
This is not an unrelated representation. For \(G=\operatorname{Sp}(6,\mathbb R)\), \(\Theta=\{\alpha_3\}\), the semisimple Levi action on \(\mathfrak u_{\alpha_3}\) is precisely this symmetric-square action. The paper describes that matrix model on pp. 723 and 732 (PDF pp. 17 and 26). The printed all-weights definition would reject a positive structure explicitly included in Theorem 3.4.\par
\subsection*{The interior need not consist of maximal-compact-fixed vectors}
Take \(S=\operatorname{SL}_2(\mathbb R)\) acting on binary quartics \(V=\operatorname{Sym}^4(\mathbb R^2)\), and set \[C_{\max}=\{f:f(x,y)\geq0\text{ for all }(x,y)\in\mathbb R^2\}.\] This is a nonzero closed pointed invariant convex cone with nonempty interior. The five distinct weight lines are connected by the raising and lowering operators, which proves irreducibility. The representation has weights \(4,2,0,-2,-4\), and is proximal. It satisfies even the printed all-weights parity condition.\par
Equation reference: eq:stab. The quartic \(q=x^4+y^4\) lies in \(\operatorname{int}C_{\max}\), because \(q\geq(x^2+y^2)^2/2\); sufficiently small coefficient perturbations stay positive on the unit circle. Its infinitesimal stabilizer satisfies, up to the harmless choice of contragredient convention, \[
(a x\partial_x+b y\partial_x+c x\partial_y-a y\partial_y)q
=4a(x^4-y^4)+4b x^3y+4c xy^3 .\] The three polynomials on the right are independent. Thus the stabilizer Lie algebra is zero and cannot contain a conjugate of \(\mathfrak{so}_2\), the one-dimensional maximal compact Lie algebra. This disproves Proposition 3.6(2)’s asserted equality.\par
General uniqueness up to sign also fails. Let \(C_{\min}\) be the cone generated by \((a x+b y)^4\). Normalize \(a^2+b^2=1\). The generators are compact and their integrals over the unit circle are a fixed positive constant; their convex hull is compact in that affine hyperplane, proving closedness of the generated cone. Five distinct slopes give independent quartics by the Vandermonde determinant, so the cone has nonempty interior. It is pointed, invariant, and contained in \(C_{\max}\). However, \[(x^2-y^2)^2\in C_{\max}\setminus C_{\min},\] since its \(x^2y^2\) coefficient is \(-2\), whereas any positive sum of fourth powers has that coefficient nonnegative. The two cones are distinct and are not negatives of one another.\par
\subsection*{Primary source, counterevidence, and repair}
Benoist’s Proposition 4.7 and proof in [B] were actually read on manuscript pp. 27–28 (PDF pp. 29–30). They use the highest restricted weight in \(2P\), together with proximality. They neither impose \(R(V)\subset2P\) nor say that all interior vectors have maximal compact stabilizers. The cited primary source therefore does not support the stronger printed claims.\par
There is substantial counterevidence against a broader negative conclusion. The classification calculation on p. 727 actually tests the extremal weight \(\beta(H_\alpha)\), consistent with the correct criterion. The Section 5 modules are special symmetric-cone models: positive matrices, Lorentz cones, and the exceptional Jordan cone. The binary quartic example is not one of these modules and does not refute their homogeneity or uniqueness.\par
The required repair is to use the usual highest-weight order and the highest-restricted-weight parity condition, retain only the valid direction that a nonzero maximal-compact-fixed vector lies up to sign in the interior, and justify uniqueness and stabilizers for the specific Section 5 models. For the one-way assertion, average an interior point over a maximal compact subgroup and use the one-dimensional fixed-vector space; this does not imply uniqueness of arbitrary invariant cones. The false general criterion and false general cone consequence share the same central dependency and are counted together as M01. The actual classified models and the repaired extremal-weight calculation must be kept in the assessment; no counterexample to the four-family list is established.\par
\section*{M02—Long restricted roots need not have multiplicity one (Minor)}
Location: the last root-multiplicity paragraph of Section 2.2, printed p. 716 (PDF p. 10). Under the general real semisimple setup it asserts that a long restricted-root space is one-dimensional.\par
Let \(G=\operatorname{Sp}(4,\mathbb C)\), viewed as a real Lie group. It is connected and real semisimple with finite center. With \[\mathfrak a=\{\operatorname{diag}(h_1,h_2,-h_1,-h_2):h_i\in\mathbb R\},\] the restricted roots form \(C_2\), with long simple root \(2e_2\), and \[\mathfrak g_{2e_2}=\mathbb CE_{24}=\operatorname{span}_{\mathbb R}\{E_{24},iE_{24}\}.\] For \(J=\left(\begin{smallmatrix}0&I\\-I&0\end{smallmatrix}\right)\), \(E_{24}^{\mathsf T}J+JE_{24}=0\) and \([H,E_{24}]=2h_2E_{24}\). The same identities hold for \(iE_{24}\); these are real-linearly independent. The real multiplicity is \(2\).\par
This refutes the unrestricted background assertion. It does not refute the classification, whose p. 727 proof separately requires multiplicity one by proximality and consults the restricted-root classification. This group fails that test. Delete the general sentence or restrict it to the families for which multiplicity one is established. The demonstrated effect is local.\par
\section*{M03—The BCH bound is local, not globally quadratic (Minor)}
Equation reference: eq:bad. Location: Lemma 10.21(2), printed pp. 768–769 (PDF pp. 62–63). It quantifies over every \(v=(v_i)\) in the indicated root spaces and asserts \[
 \left\|\log(\exp v_1\cdots\exp v_N)-\sum_i v_i\right\|
 \leq C\max_i\|v_i\|^2 .\] No smallness restriction is imposed.\par
Equation reference: eq:bch. For split \(\operatorname{SL}_4(\mathbb R)\), use the longest reduced word \((1,2,3,1,2,1)\), put \(v_1=tE_{12},v_2=tE_{23},v_3=tE_{34}\), and set the other three parameters to zero, with \(t>0\). These lie even in the closed positive cones. Multiplication and the terminating formula \(\log(I+N)=N-N^2/2+N^3/3\) give \[
 \log(\exp v_1\exp v_2\exp v_3)
 =t(E_{12}+E_{23}+E_{34})
  +\frac{t^2}{2}(E_{13}+E_{24})+\frac{t^3}{3}E_{14}.\] The squared Frobenius norm of the remainder is \(t^4/2+t^6/9\), so its ratio to \(t^2\) is unbounded. This contradicts [eq:bad] (\#eq:bad).\par
There is a complete repair sufficient for the downstream use. The finite BCH polynomial gives, for \(r=\max_i\|v_i\|\leq1\), \[\|R(v)\|\leq\sum_{k=2}^{d}C_kr^k\leq
 \left(\sum_{k=2}^{d}C_k\right)r^2 .\] Proposition 10.10 bounds \(r\leq D\|X\|\) for \(X=\log F_\gamma(v)\) in the closed cone parameterization. This proves Proposition 10.22(1) near zero. Away from zero, say \(\|X\|\geq\delta\), use \[\|X-p_1X\|\leq\|I-p_1\|\|X\|
 \leq\frac{\|I-p_1\|}{\delta}\|X\|^2.\] Thus the conclusion of Proposition 10.22(1) still has a global proof. Its tangent-cone limit and Trotter product argument need only the local BCH estimate. The unbounded-parameter counterexample does not invalidate those conclusions or Theorem 10.24. This bounded impact is the reason for Minor severity.\par
\section*{Additional mathematical checks}
The reviewed semigroup proof uses injectivity through Bruhat lengths, openness through invariance of domain, properness through the abelianization, exclusion of boundary parameters from the open Bruhat cell, and a common split-subgroup component to compare reduced words (Sections 10.2–10.7). The properness inequality has a direct proof: choose a linear functional positive on the compact unit section of a pointed closed cone; applying it to a sum bounds each summand by the norm of that sum. This justifies the substantive estimate in Proposition 10.10. The argument first gives a homeomorphism; a complete differential-rank proof for every claimed diffeomorphism was not reconstructed.\par
The open/closed distinction is essential. Remark 10.20 and Section 11.2 explicitly allow noninjectivity on the entire closed parameter cone. Vanishing scalar parameters merge vector parameters. This is acknowledged behavior, not a new finding.\par
For Lemma 11.2, write \(A=q(v_1)\), \(B=b(v_1,v_2)\), \(C=q(v_2)\), and \(S=q(v_1+v_2)=A+2B+C\). Then \[q\big(Sv_1-A(v_1+v_2)\big)
 =S^2A-2SA(A+B)+A^2S=SAC.\] For future timelike vectors this is positive; connectedness and \(v_1=v_2\) fix the future component. All seven transition equations (11.2)–(11.8) also passed an exact rational example with strictly interior inputs and outputs.\par
The Appendix C reflection matrices were reconstructed from its four \(F_4\) roots. The product \(\sigma_1\sigma_2\) has order six and \((\sigma_1\sigma_2)^3s_3s_4s_3=-I_4\). The root system has 24 positive roots, matching the expanded word length. This verifies the displayed longest-element calculation, not all earlier exceptional Lie-algebra constructions.\par
For the \(\operatorname{SL}_3\) examples of Section 12.4, multiplication yields \[r(a+a')(b+b')-(c+c'+ab')
 =(rab-c)+(ra'b'-c')+ra'b+(r-1)ab'\geq0\] when \(r\geq1\). Larger sharp invariant semigroups therefore really occur; Theorem 12.1 concludes inclusion rather than equality. The stronger Bruhat-component assumption in Theorem 12.2 must be preserved.\par
For flags, uniqueness of the opposite diamond follows from inversion of the same semigroup in a fixed unipotent chart. Transversality between opposite diamonds reduces to multiplication of positive elements. The distinction between \(\operatorname{Aut}_1(\mathfrak g)\) and \(G\), and the finite-index qualification in the properness proof for positive triples, were preserved.\par
\section*{Coverage and limitations}
Text was actually read on PDF pp. 1–10, 16–29, 55–85, and 90–93: 59 distinct pages. Twenty-four source pages were visually inspected, including all formal finding locations. Full extraction and indexing are not counted as scientific verification.\par
This claim-focused initial review is not a correctness proof of all 93 pages. The entire \(\Theta\)-system and split-subalgebra construction in Sections 5.5–7, all Bruhat-length inputs in Sections 8–9, the case-by-case boundary parameterization of Proposition 11.5, the bracket-generation argument in Lemma 12.9, the continuity step in Proposition 13.4, and the remaining cyclic-order discussion were not fully reconstructed. Some downstream checks are conditional on these inputs. Exceptional Jordan theory and the complete external bibliography were not reproved.\par
Local sign, index, and base-case slips with an immediate intended correction were separated from substantive findings. For example, Proposition 13.26’s quadruple criterion needs \(n\geq4\); for \(n=3\) the positive-triple condition must remain because the displayed quadruple test is empty. The contraction parameter in Proposition 13.4 must likewise have the sign consistent with contraction of positive root spaces. Neither observation is used to allege a false main positivity theorem.\par
The acquired manifest has the publisher PDF and metadata, with no separately bound author program or supplement. A limited version/correction search and earlier author-preprint metadata do not establish absence of all corrections. No current erratum repairing these specific statements was established. Independent specialist examination of the repairs and human acceptance remain pending.\par
9 O. Guichard and A. Wienhard, Generalizing Lusztig’s total positivity, Invent. Math. 239 (2025), 707–799. DOI: 10.1007/s00222-024-01303-y. Assigned publisher version.\par
Y. Benoist, Automorphismes des cônes convexes, Invent. Math. 141 (2000), 149–193. DOI: 10.1007/PL00005789. Author manuscript, Proposition 4.7 and proof, manuscript pp. 27–28 (PDF pp. 29–30), actually read from https://www.imo.universite-paris-saclay.fr/\textasciitilde{}yves.benoist/prepubli/0099automo.pdf. Manuscript pagination is not asserted to equal journal pagination.\par
\end{document}
